\begin{textsample}{POS Dim 3 – ma – Score 9.00 – ma/ma\_inf\_7.txt}  \label{ex:f3_pos_167}
5.
Espaços, tempos, quantidades, relações e transformações O campo de experiências “ espaços, tempos, quantidades, relações e transformações ” integra vivências que proporcionam à \textbf{criança}, na sua relação com o meio ambiente, investigar, questionar, comunicar quantidades, explorar o espaço e os objetos, estabelecendo relações entre eles, transformando -os e ressignificando -os, a partir das \textbf{brincadeiras}, das interações e do estímulo com materiais e espaços variados.
Nos termos da BNCC : > As \textbf{crianças} vivem inseridas em espaços e tempos de diferentes dimensões, em um mundo constituído de fenômenos naturais e socioculturais.
Desde muito \textbf{pequenas}, elas procuram se situar em diversos espaços ( rua, bairro, cidade etc. ) e tempos ( dia e noite ; hoje, ontem e amanhã etc. ). \textbf{Demonstram} também \textbf{curiosidade} sobre o mundo físico ( seu próprio corpo, os fenômenos atmosféricos, os animais, as plantas, as transformações da natureza, os diferentes tipos de materiais e as possibilidades de sua manipulação etc. ) e o mundo sociocultural ( as relações de parentesco e sociais entre as pessoas que conhece ; como vivem e em que trabalham essas pessoas ; quais suas tradições e seus costumes ; a diversidade entre elas etc. ).
Além disso, nessas experiências e em muitas outras, as \textbf{crianças} também se deparam, \textbf{frequentemente}, com conhecimentos matemáticos ( contagem, ordenação, relações entre quantidades, dimensões, medidas, comparação de \textbf{pesos} e de comprimentos, avaliação de distâncias, reconhecimento de formas geométricas, conhecimento e reconhecimento de numerais cardinais e ordinais etc. ) que igualmente aguçam a \textbf{curiosidade}.
Portanto, a Educação \textbf{Infantil} precisa promover experiências nas quais as \textbf{crianças} possam fazer observações, \textbf{manipular} objetos, investigar e explorar seu entorno, levantar hipóteses e consultar fontes de informação para buscar respostas às suas \textbf{curiosidades} e indagações.
Assim, a instituição escolar está criando oportunidades para que as \textbf{crianças} ampliem seus conhecimentos do mundo físico e sociocultural e possam utilizá -los em seu cotidiano ( BRASIL, 2017:40 ).
Por meio de práticas cotidianas permeadas de situações significativas e estruturadas de experiências em que as \textbf{crianças} são protagonistas, elas têm oportunidade de quantificar, medir, formular hipóteses, solucionar problemas, comparar e orientar -se no espaço e no tempo, com ricas possibilidades de conexão com o aparato científico e tecnológico, além de aprender a valorizar a vida no planeta.

% matched lemmas: brincadeira, criança, curiosidade, demonstrar, frequentemente, infantil, manipular, pequeno, peso
\end{textsample}
